
This work investigates whether the per-sample normalized last-layer gradient norm $\tilde{g}_i = \|\nabla _W \ell _{i}\| / \|h(x_{i})\|$ provides a usable minority group proxy signal during standard ERM training, without group annotations. On Waterbirds, $\tilde{g}$ achieves AUC = 0.914 for minority group prediction at epoch 5 of ERM training, with a minority/majority ratio of 8.8x. The signal persists to epoch 10 (ratio 8.5x, AUC 0.888). Feature norm equalization (h\_norm\_std\_ratio $\approx$ 0.10) is confirmed, supporting the outer-product decomposition interpretation that the gradient norm disparity reflects prediction residual differences rather than feature-scale artifacts.

The existence experiment (H-E1) receives a PARTIAL gate result (2/3 criteria pass). The third criterion --- balance deviation of the top-25\% subset --- fails at 0.379 against a $\leq$ 0.10 target; post-hoc analysis identifies this as a criterion design mismatch rather than a mechanism failure. A revised experiment with minority recall as the criterion was proposed but not executed.

The full GNR-LLR pipeline (Stage 2 last-layer retraining) was not implemented or evaluated. No WGA results exist. The relationship between the strong proxy signal quality reported here and downstream WGA improvement via Stage 2 retraining remains to be established empirically.

Three directions follow from this work: (1) executing the full GNR-LLR pipeline (Stage 2) to measure WGA on Waterbirds; (2) running h-e1-v2 with minority recall as the criterion and comparing AUC($\tilde{g})$ against AUC(binary misclassification); (3) evaluating the proxy signal on CelebA.

